\documentclass[10pt,a4paper]{amsart}
\usepackage[margin=19mm]{geometry}
\usepackage{amsmath,amssymb,mathtools}
\usepackage[scr=rsfs,cal=euler]{mathalfa}
\usepackage{xfrac}
\usepackage[unicode,hidelinks,bookmarks=false]{hyperref}
\newcommand{\ie}{i.e.\ }
\newcommand{\cf}{cf.\ }
\newcommand{\resp}{respectively\ }
\newcommand{\Subsection}{\subsection}
\providecommand{\nobreakdash}{\nobreak\hbox{-}}
\setlength{\emergencystretch}{2em}
\multlinegap0pt
\usepackage{bm}
\usepackage{bbm}
\usepackage{dsfont}
\usepackage{xspace}
\usepackage{cancel}
\usepackage{mathtools}
\usepackage{amsthm}
\makeatletter
\@ifundefined{theorem}{\newtheorem{theorem}{Theorem}}{}
\@ifundefined{claim}{\newtheorem{claim}[theorem]{Claim}}{}
\@ifundefined{lemma}{\newtheorem{lemma}[theorem]{Lemma}}{}
\makeatother
\newcommand{\R}{\mathbb{R}}
\newcommand{\asc}{\mathscr{A}}
\newcommand{\des}{\mathscr{D}}
\newcommand{\g}{\mathfrak{g}}
\title[Equivariant Morse-Bott cohomology through stabilization]{Equivariant Morse-Bott cohomology through stabilization}
\author{Erkao Bao, Robi Huq, Shengzhen Ning}
\date{Source: arXiv:2601.16119v1 (CC BY 4.0)}
\begin{document}
\maketitle

\noindent\emph{Source and licence.} Equivariant Morse-Bott cohomology through stabilization. Version arXiv:2601.16119v1 (CC BY 4.0), distributed under the Creative Commons Attribution 4.0 International licence, as recorded in the arXiv licence metadata for this exact version and shown on the abstract page https://arxiv.org/abs/2601.16119v1. Full rights notice: licenses/arxiv-2601.16119-CC-BY-4.0.txt.\par\smallskip

\noindent\textit{Editorial scope.} Counted unit: the Lemma whose source label is \texttt{lem:equivMBS} in the article (source0.tex lines 741--743, source label \texttt{lem:equivMBS}), delivered together with the article's own complete original proof of it (source0.tex lines 744--765, one \texttt{proof} environment of 22 lines). No mathematics is written, repaired or re-derived: statement and proof are the article's own text line for line. Required context retained uncounted: none. Every definition, display and label of those lines is retained unchanged. Omitted from this package and not redistributed: 1--740, 766--1529. Nothing inside a retained excerpt is omitted or altered: every retained body is delivered exactly as the article's lines are written, so no omission receipt of any kind accompanies this package. Citations: the retained excerpts carry no citation command at all, so no bibliography entry of the article is transcribed and no reference is redistributed. Reference adaptation: none was needed and none was made; every reference inside the delivered unit resolves inside it, so no unresolved reference or citation is produced. Printed slips kept literal: no printed word, symbol, hypothesis or number of the counted statement or of its proof was altered. Layout and class: the article's own class is \texttt{amsart} with packages including amsmath, amssymb, amsbsy, amscd, mathtools, times, mathrsfs, euscript, thmtools, epsf, overpic, psfrag, stmaryrd, caption; the wrapper is the lane's pinned \texttt{amsart} editorial wrapper at 10pt on A4, which restates the theorem-like environments the retained text needs and adds the article's own macro definitions that the retained text needs (\texttt{\textbackslash R}, \texttt{\textbackslash asc}, \texttt{\textbackslash des}, \texttt{\textbackslash g}), with extra wrapper packages (bm, bbm, dsfont, xspace, cancel, mathtools). The delivered numbering is the unit's own. Assets: no retained span contains a figure environment, a graphics-inclusion command, a PDF-page inclusion, a table or a TikZ picture; no image file is copied into this package, the package declares no asset dependency and no image file is redistributed with it. Licence: version arXiv:2601.16119v1 of this article is distributed under the Creative Commons Attribution 4.0 International licence, as recorded in the arXiv licence metadata for this exact version and shown on the abstract page https://arxiv.org/abs/2601.16119v1. A full rights notice is delivered at licenses/arxiv-2601.16119-CC-BY-4.0.txt. Characters: every retained excerpt is pure ASCII, so the delivered engine, XeLaTeX with its default Latin Modern text font, reports no missing character.
\begin{lemma}[Weakly MBS implies MBS]\label{lem:equivMBS}
    Let the compact Lie group $G$ act on the closed manifold $M$ with a $G$-Morse-Bott function $f$ and a $G$-invariant metric $g$. Assume $G\cdot p$ and $G\cdot q$ are two critical orbits of $f$. Then $\des(p)\pitchfork_M \asc(G\cdot q)$ if and only if $\des(G\cdot p)\pitchfork_M \asc(G\cdot q)$.
\end{lemma}

\begin{proof}
    We only need to prove the `if' direction. For any $x\in \des(p)\cap\asc(G\cdot q)$, assume \[T_x\des(G\cdot p)+T_x\asc(G\cdot q)=T_xM.\] Our goal is to show \[T_x\des(p)+T_x\asc(G\cdot q)=T_xM.\] We may further assume that $x$ is arbitrarily close to the critical point $p$ due to the following simple fact.
    \begin{claim}\label{claim:transverse}
   Let $\gamma: \R \to M$ be a gradient flow line of $(f,g)$ from $p$ to some point in $G\cdot q$. Then for any $s, s' \in \R$, $\des_{g}(p)$ and $\asc_{g}(G\cdot q)$ intersect transversely in $M$ at $\gamma(s)$ if and only if they  intersect transversely at $\gamma(s')$.
\end{claim}
\begin{proof}[Proof of Claim]
This follows from the observation that the descending and ascending manifolds are invariant under the gradient flows. So, it maps $T_{\gamma(s)} \des(p)$ isomorphically to $T_{\gamma(s')} \des(p)$, and maps $T_{\gamma(s)} \asc(G\cdot q)$ isomorphically to $T_{\gamma(s')} \asc(G\cdot q)$. 
\end{proof}

    
    On the one hand, since the pair $(f,g)$ is $G$-invariant, the orbit $G\cdot x\subseteq\des(G\cdot p)$ which implies that 
    \[T_x\des(p)+T_x(G\cdot x)\subseteq T_x\des(G\cdot p).\]
    
    On the other hand, let us take a basis in the Lie algebra  $\g$ and consider their corresponding fundamental vector fields $X_1,\cdots,X_{\text{dim}G}$ on $M$ which span the tangent spaces of the orbits. At the point $p$, we have\[0=T_p\des(p)\cap T_p(G\cdot p)=T_p\des(p)\cap \text{span}\{X_1,\cdots,X_{\text{dim}G}\}|_p.\] For dimensional reason, it follows that \[T_p\des(G\cdot p)=T_p\des(p)+\text{span}\{X_1,\cdots,X_{\text{dim}G}\}|_p.\] Since $x$ is sufficiently close to $p$, by continuity we have \[T_x\des(G\cdot p)=T_x\des(p)+\text{span}\{X_1,\cdots,X_{\text{dim}G}\}|_x=T_x\des(p)+ T_x (G\cdot x).\] 
    
    Finally, the $G$-invariance of $(f,g)$ also implies that $T_x(G\cdot x)\subseteq T_x\asc(G\cdot q)$. Therefore, we can see that 
    \begin{align*}
        T_x\des(p)+T_x\asc(G\cdot q)&= T_x\des(p)+T_x(G\cdot x)+T_x\asc(G\cdot q)\\
        &=T_x\des(G\cdot p)+T_x\asc(G\cdot q)\\
        &=T_xM.
    \end{align*}
\end{proof}

\end{document}
